\documentclass[11pt]{article}
\usepackage[margin=1in]{geometry}
\usepackage{hyperref}
\usepackage{booktabs}
\usepackage{amsmath,amssymb}
\usepackage{listings}

\title{Replication Report: Boixo et al.\ (2017)\\
``Simulation of low-depth quantum circuits as complex undirected graphical models''}
\author{Ollie (OpenClaw AI subagent) --- QC-100 Replication Wave}
\date{2026-07-03 (backfilled 2026-07-06)}

\begin{document}
\maketitle

\begin{abstract}
We independently replicate the Boixo, Isakov, Smelyanskiy \& Neven
(arXiv:1712.05384) proposal to simulate shallow universal quantum circuits by
representing the amplitude $\langle x|U|0\rangle$ as an undirected graphical
model with complex factors, evaluated via variable elimination (equivalently,
tensor--network contraction). On a sweep of 70 random shallow circuits
($n=8$--$16$, depth $d=2$--$6$, geometries $1{\times}m$ / $2{\times}m$ /
$3{\times}m$ / $4{\times}4$) run on a single laptop CPU with \texttt{numpy} +
\texttt{opt\_einsum} (greedy path), we confirm four testable claims of the
paper: (C1) exact equivalence of tensor-network and Schr\"odinger amplitudes
($\max|\Delta| = 3.4{\times}10^{-17}$); (C2) the contraction width is bounded
by $\min(O(d\cdot\ell), O(n))$ with an $O(1)$ constant; (C3) per-depth
$\log_2$FLOPs slope $\approx \ell_{\min}$; (C4) TN-cost / $2^n$ falls below 1
for $n\gtrsim 10$ at $d=2$, reaching $\sim\!2.4\times 10^{-2}$ at $n=16$.
\textbf{Verdict: REPLICATED} on the functional-form / scaling claims. The
paper's supercomputer-scale numerical points (e.g.\ $7{\times}8$ depth 30) are
out of scope for a small-instance replication and were not attempted.
\end{abstract}

\section{Paper summary}
The paper introduces a classical simulator for shallow universal random
quantum circuits based on mapping the amplitude $\langle x|U|0\rangle$ to an
\emph{undirected graphical model with complex factors} (equivalently a tensor
network), and evaluating it via variable elimination. The runtime is
exponential in the \emph{treewidth} of the induced graph, which for 2D
nearest-neighbor circuits scales as $\min(O(d\cdot\ell), O(n))$, where $d$ is
depth, $\ell$ is the smaller lateral dimension of the qubit lattice, and $n$
the total qubit count. On a single 128\,GB workstation with QuickBB
elimination-ordering search plus TensorFlow/Dask contraction, the authors
compute amplitudes for circuits up to $5{\times}9$ depth 40, $7{\times}8$
depth 30, and $10{\times}\kappa$ depth 19, previously done only on Cori II /
Blue Gene/Q.

\section{Claims tested}
\begin{center}\small
\begin{tabular}{clll}
\toprule
\# & Claim & Type & Tested? \\
\midrule
C1 & TN $\equiv$ Schr\"odinger amplitude & Correctness & Yes (70/70) \\
C2 & width $\le \min(O(d\ell), O(n))$ & Scaling & Yes (70/70) \\
C3 & FLOPs exponential in treewidth, slope $\approx \ell_{\min}$ & Scaling & Yes, per grid \\
C4 & TN cost $<2^n$ for shallow $d$ & Scaling & Yes, ratio down to $2.4{\times}10^{-2}$ \\
C5 & Supercomputer-scale $7{\times}8$ depth 30 in Fig.~3 & Numerical @ scale & \textbf{No} (out of scope) \\
\bottomrule
\end{tabular}
\end{center}

\section{Method (summary)}
Real numerical simulation, no fabricated numbers, fresh Python
3.14.6 / \texttt{numpy} 2.5.0 / \texttt{opt\_einsum} 3.4.0 venv on CherryRd
(macOS, laptop CPU). Circuit generator: Hadamard layer, then $d$ layers of
random single-qubit gates from $\{T, \sqrt{X}, \sqrt{Y}\}$ (paper's menu) on
qubits touched by the prior 2-qubit layer, followed by nearest-neighbor CZs
on a periodic \{H-even, H-odd, V-even, V-odd\} pattern (structurally faithful
to the paper's supremacy-style circuits). Ground truth: full-statevector
Schr\"odinger evolution over $\mathbb{C}^{2^n}$. TN builder: per-qubit
increasing wire indices; single-qubit gates $\to$ rank-2 tensors; CZs $\to$
rank-4 diagonal tensors; close each wire with $|0\rangle$; contract to scalar
via \texttt{opt\_einsum} greedy path. Cost extracted from \texttt{PathInfo}:
\texttt{largest\_intermediate}, \texttt{contraction\_width} =
$\log_2(\texttt{largest\_intermediate})$, \texttt{opt\_cost\_flops}. Full
reproducer commands are in \texttt{report/REPORT.md} \S3f and
\texttt{workflow.md}.

\section{Results vs paper (headline numbers)}
\paragraph{C1 (correctness).} $\max|\text{TN}-\text{SV}| = 3.42\times 10^{-17}$
across all 70 configs. \textbf{REPLICATED}.

\paragraph{C2 (width bound).} $\mathrm{width}/\min(d\ell, n)$: min $0.333$,
median $0.667$, mean $0.810$, max $2.000$. All ratios $O(1)$. \textbf{REPLICATED}.

\paragraph{C3 (slope in depth).} Log-linear fit
$\log_2(\text{FLOPs}) = a\cdot d + b$ per grid gives $a\!\approx\!0.4$ at
$\ell_{\min}=1$, $0.7$--$1.0$ at $\ell_{\min}=2,3$, and $\approx 1.0$ at
$\ell_{\min}=4$; monotonically increasing with $\ell_{\min}$. \textbf{REPLICATED}
(with $O(1)$ prefactor smaller than the paper's vertical-ordering baseline;
greedy \texttt{opt\_einsum} beats vertical on small instances, consistent
with the paper's remark that QuickBB improves on vertical).

\paragraph{C4 (crossover).} At $d=2$, $\text{TN FLOPs}/2^n$ falls from
$3.10$ ($n{=}8$) to $2.4\times 10^{-2}$ ($n{=}16$). TN wins wall-clock in
$38/70$ configs already at these tiny sizes, and dominates on every
$n{\ge}14$, $d{\le}5$ point. \textbf{REPLICATED}.

\section{Critical assessment (independent replication caveats)}
This backfill applies the hard requirement (Rick, 2026-07-05) that every QC
replication carry a genuine critique. Key limitations of the present
replication:

\begin{enumerate}
\item \textbf{Graphical-model mapping was verified structurally + numerically,
but not against the paper's own reference implementation.} We built the TN
mapping from the description in the paper and verified TN amplitudes match
statevector amplitudes to machine precision. This confirms our mapping is
\emph{some} valid TN for the circuit --- it does not certify it is
byte-identical to the mapping the authors used. A stronger test would be to
run the paper's own code (Google / QuickBB pipeline) on identical circuit
inputs and diff FLOP counts + contraction widths, which we did not do.
\item \textbf{Hardness / simulability claims were not reproduced at the
paper's specific circuits.} The paper's headline empirical points
($5{\times}9$ depth 40, $7{\times}8$ depth 30, $10{\times}\kappa$ depth 19)
were not attempted because the QuickBB elimination-ordering search alone
takes $\sim$1 day on those sizes and the contraction needs $\ge$128 GB RAM.
The scaling exponents are reproduced on a $n{\le}16$, $d{\le}6$ sweep, and
the exponents are then \emph{extrapolated} to argue the paper's supercomputer
points are consistent --- this is a functional-form verification, not a
number-for-number verification.
\item \textbf{No independent comparison against alternative low-depth
simulators was made.} Modern classical-simulation benchmarks
(tensor-network methods like \texttt{quimb} / \texttt{ITensor} /
\texttt{cotengra}, PEPS approximations, sparse-simulator Feynman-path methods
like \texttt{qFlex}, and hybrid Schr\"odinger--Feynman approaches) all
compete on the same regime. We only compare against direct statevector, so
we cannot say whether the paper's approach is state-of-the-art among modern
graphical-model / tensor-network simulators.
\item \textbf{Elimination-ordering heuristic differs.} The paper uses QuickBB
(SAT-style branch-and-bound over elimination orderings, potentially better
than greedy for large graphs). We use \texttt{opt\_einsum}'s greedy heuristic
(a linear-time local heuristic). For small instances the two are comparable
(both find good orderings), but at the paper's scale QuickBB is known to
find tighter orderings, so our small-instance $O(1)$ prefactor should not be
extrapolated to the paper's regime as a competitive claim.
\item \textbf{Circuit generator matches the paper's menu (H, T, $\sqrt{X}$,
$\sqrt{Y}$, CZ) but not the exact random seeds.} We generated new circuits
with our own RNG; we cannot re-derive the paper's specific per-instance
numbers. This is the standard replication-vs-reproduction distinction and is
appropriate for a scaling-claim paper, but it does mean we cannot spot-check
individual amplitudes in Figures 3--5 against ours.
\item \textbf{No noisy-circuit extension attempted.} The paper is purely
noiseless. A live open question (see \S6) is whether the graphical-model
representation extends cleanly to noisy channels; we did not attempt this.
\end{enumerate}

Bottom line: the mapping and the scaling machinery replicate cleanly at
laptop scale; the strongest claims of the paper (the specific
$5{\times}9$/$7{\times}8$/$10{\times}\kappa$ numerical results, and any
implicit ``QuickBB + TensorFlow beats all comers'' state-of-the-art claim)
were not tested here.

\section{Open questions}
See \texttt{report/open\_questions.json} (5 concrete probes with next steps).

\section{Verdict}
\textbf{REPLICATED} on the functional-form / scaling claims (C1--C4).
Out-of-scope for the supercomputer-scale numerical points (C5). See
\S5 for critique-scoped caveats.

\section{Artifacts}
See \texttt{report/artifacts\_summary.md} for the file-level manifest, and
\texttt{report/workflow.md} for the end-to-end reproduce recipe. Full
prose report is in \texttt{report/REPORT.md}.

\end{document}
